\chapter{讨论与结论}
\label{ch:discuss}

\section{结果要点}

\begin{enumerate}[leftmargin=2em]
  \item \textbf{可部署性成立.}
  $z_{\mathrm{rob}}(\hat H)$+Adaptive-MKL 不依赖真 $H$ 与真后验 $f^*$；Oracle（真 $H$+$f^*$）与 MLD 一致，说明在结构坐标上核逼近有效，主线差距主要来自 CSI/坐标估计而非“核方法无效”。
  \item \textbf{中期线性区有稳定正增益.}
  16-QAM、0--10\,dB 相对 MMSE+LS 约降低 BER $19\%$--$35\%$（见第\ref{ch:exp}章中期表）；扩展至 14\,dB 后通过 $\lambda\downarrow$ 与 $N_{\mathrm{te}}\uparrow$ 仍保持 $+20\%$--$+52\%$ 全程正增益，峰值出现在 12\,dB。
  \item \textbf{前端失配是优势工况.}
  同一流水线在 soft-clip / Kerr / MZM / hard-clip 等幅度型失真下，平均 $G$ 约 $23\%$--$52\%$；软限幅尤为突出。原因是线性 MMSE 仍按 $Y=HX+N$ 工作，而 RKHS 在 $z_{\mathrm{rob}}$ 上直接拟合失配后的判决边界。
  \item \textbf{更真实信道可迁移.}
  Kronecker 相关 + soft-clip 平均 $G\approx 39\%$（12\,dB 约 $61\%$）；在为每用户打散到达角后的 CDL-A 上，soft-clip 平均约 $19\%$（表\ref{tab:kron-gain}--\ref{tab:cdl-gain}）。这说明收益不只建立在 i.i.d.\ Rayleigh 上。
  \item \textbf{64-QAM 可直接切换星座.}
  绝对 BER 升高符合星座加密；相对 MMSE 的增益模式与 16-QAM 一致。16-QAM 线性经 $\lambda$ 降至 $0.01$--$0.03$ 与 $N_{\mathrm{te}}=2\times10^{4}$ 后，0--14\,dB 全程正增益；64-QAM 线性在 16--20\,dB 仍可能出现负 $G$，属稀有错误 + 有限样本核估计的已知瓶颈。
  \item \textbf{soft-clip 下的 MLD“地板”不是实现错误.}
  前端压缩后，按线性模型计分的 MLD/Oracle 会出现约 $0.27$--$0.32$ 的误码地板；可部署 RKHS 仍可相对 MMSE 拉开差距。解读时勿把该地板误当成信道病态。
\end{enumerate}

\section{工程取舍}

\begin{itemize}[leftmargin=2em]
  \item 坚持可部署约束：不用真 $H$、不用 $f^*$ 监督；Oracle 仅作上界对照。
  \item NN 只生成核系数，不替换 $K_\eta\alpha^\top$ 决策形式。
  \item 高 SNR：减弱稳健加载、关闭堆叠、减小正则；已赢 MMSE 或处于高误码地板时不再空转 $\lambda$ 重试。
  \item CDL 多用户必须角度可分；禁止把单链路 CDL 原样复制 $K$ 次。
\end{itemize}

\section{后续工作}

\begin{enumerate}[leftmargin=2em]
  \item 高 SNR 与线性 MMSE 的融合/切换；
  \item 更大规模系统级信道、导频污染与导频开销敏感性；
  \item Nyström / 随机特征以降复杂度，并给出相对 MMSE 的时延对照；
  \item 由只检 $X_1$ 扩展到多用户软信息输出。
\end{enumerate}

\section{结论}

本报告给出一套以 RKHS 为假设类、以贝叶斯后验损失为目标的大规模 MU-MIMO 单用户检测方案。通过 $z_{\mathrm{rob}}$、Adaptive-MKL 与核内增强，在 i.i.d.\ / Kronecker / CDL、多种前端失真及 16/64-QAM 设定下相对 MMSE+LS 取得可量化的 BER 增益；方法形式保持核展开判决，满足可部署约束，具备工程对照与继续产品化的基础。
